%%%%%%%%%%%%%%%%%%%%%part.tex%%%%%%%%%%%%%%%%%%%%%%%%%%%%%%%%%%
% 
% sample part title
%
% Use this file as a template for your own input.
%
%%%%%%%%%%%%%%%%%%%%%%%% Springer Nature %%%%%%%%%%%%%%%%%%%%%%

\begin{partbacktext}
\part{Theory}
\noindent{This part of the book is entirely dedicated to the formulation
of a theoretical basis and terminology used throughout the book. It therefore
relies on necessary, heavy mathematics and rather complex graphical 
illustrations to ensure that every piece of vocabulary is distinguishable 
by its definition. After reading this part thoroughly and unterstanding 
its content, the reader is ready to solve the problems provided at 
the end of each chapter. To give a brief overview, the topics included
are first of all some mathematical notation and definitions of vital 
preliminaries. Furthermore, the basics of rendering shall
be explained and some integration methods will be covered.}
\end{partbacktext}
